\providecommand{\BM}{\textbf{[B]}}
\providecommand{\KM}{\textbf{[K]}}
\providecommand{\SM}{\textbf{[S]}}
\providecommand{\EM}{\textbf{[E]}}
\providecommand{\XM}{\textbf{[X]}}

\sloppy
\emergencystretch=2em
\chapter*{The Matrix-Element Principle and Its Reach}
\addcontentsline{toc}{chapter}{The Matrix-Element Principle and Its Reach}

{\small\noindent\textbf{Abstract.}\quad\ignorespaces
The lepton masses follow without free parameters from a single matrix
element, $|\langle v_0|R_5|v_1\rangle|^2 = 2/9$. This document examines what
the same mechanism carries beyond that. First result \BM: in the
$\mathbb Z_3$ mode basis the icosahedral group $A_5$ generates exactly eight
squared matrix-element moduli, $\{0,\ (5-3\phi)/9,\ 1/9,\ 2/9,\ 4/9,\ 5/9,\
(2+3\phi)/9,\ 1\}$; the reach of the mechanism is therefore exactly bounded.
Second result \KM: the on-shell Weinberg angle $1-M_W^2/M_Z^2$ lies within
$0.4\,\%$ of $2/9$ --- the on-shell scheme is the natural one for pole masses
from $\tilde T\cdot m=1$; the $\overline{\rm MS}$ deviation of $3.9\,\%$ is
scheme dependence. Prediction $M_W = 80.420$\,GeV. Third result \KM: the
atmospheric mixing angle satisfies $\sin^2\theta_{23}\in\{4/9,\,5/9\}$; the
octant ambiguity of the data is the ambiguity $4/9\leftrightarrow 5/9$,
maximal mixing $1/2$ is excluded. Negative \XM: $\sin^2\theta_{12}$,
$\sin^2\theta_{13}$ and $\alpha_s$ do not lie in the eight-element set and
cannot be reached by a single $A_5$ matrix element.
}\medskip

\section*{Notation}
\addcontentsline{toc}{section}{Notation}

\medskip
\textbf{Epistemic markers:}\par\medskip
\begin{tabular}{@{}ll@{}}
\textbf{[B]} & algebraically proven (check script, exact arithmetic) \\
\textbf{[K]} & numerically confirmed \\
\textbf{[S]} & structurally plausible, identification open \\
\textbf{[E]} & established external mathematics or physics \\
\textbf{[X]} & tested and negative \\
\end{tabular}

\medskip
\begin{tabular}{@{}lp{9cm}@{}}
\toprule
\textbf{Symbol} & \textbf{Meaning} \\
\midrule
$R_5$ & five-fold rotation ($72^\circ$), order-5 element of $A_5$ \\
$v_0,v_1,v_2$ & orthonormal $\mathbb Z_3$ modes (symmetric, phase-carrying) \\
$p_k = |\langle v_0|R_5|v_k\rangle|^2$ & icosahedral redistribution weights \\
$\phi = (1+\sqrt5)/2$ & golden ratio \\
$\theta = 2/9 = p_0$ & Koide phase (= Born probability, Doc.~367) \\
$\xi = 4/30000$ & FFGFT scaling parameter \\
$\sin^2\theta_W$ & Weinberg angle (electroweak mixing) \\
$\theta_C$ & Cabibbo angle (quark generation mixing) \\
\midrule
$C_3$ & three-fold rotation ($120^\circ$), second generator of $A_5$ \\
$A_5$ & icosahedral group (60 elements) \\
$\mathbb Z_3$ & cyclic group of order 3 (mode symmetry) \\
$\mathrm{GF}(q)$, $\mathrm{GF}(q)^*$ & Galois field with $q$ elements, its multiplicative group \\
$\tilde T$ & proper-time scaling ($\tilde T\cdot m = 1$) \\
$Q$ & Koide quotient $\sum m_i/(\sum\sqrt{m_i})^2$ \\
$Q_{\rm up}$, $Q_{\rm down}$ & Koide quotient of the up and down quarks \\
$r_i$, $p_i$ & prefactor and exponent of the mass formula $m_i = r_i\,\xi^{p_i}v$ \\
$K_{\rm frak}$ & fractal correction factor $1-100\xi$ (Doc.~133) \\
\midrule
$M_W$, $M_Z$ & masses of the $W$ and $Z$ bosons (pole masses) \\
$m_h$, $m_t$, $m_s$ & mass of the Higgs particle, of the top quark, of the strange quark \\
$m_e$, $m_p$ & mass of the electron, of the proton \\
$v$ & electroweak vacuum scale ($246.22$\,GeV, SM definition) \\
$\overline{\rm MS}$ & renormalisation scheme (minimal subtraction) \\
$\theta_{12}$, $\theta_{13}$, $\theta_{23}$ & neutrino mixing angles (solar, reactor, atmospheric) \\
$\delta_{CP}$ & CP-violating phase \\
$\Delta m^2_{\rm atm}$, $\Delta m^2_{\rm sol}$ & differences of the squared neutrino masses \\
$\alpha$, $\alpha_s$ & fine-structure constant, strong coupling \\
$\Lambda_{\rm QCD}$ & confinement scale of the strong interaction \\
$R_{\rm Proton}$ & charge radius of the proton \\
$r_e$, $\lambda_p$ & classical electron radius, Compton wavelength of the proton \\
$\sigma$ & standard deviation of the measurement \\
\bottomrule
\end{tabular}

\section{The Mechanism and Its Yield for the Leptons}
\label{sec:ertrag}

The starting point of the mechanism is a single number: the squared modulus
of the matrix element of the five-fold rotation $R_5 \in A_5$ between the
symmetric mode $v_0$ and the electron mode $v_1$,
\[
  p_0 = |\langle v_0|R_5|v_1\rangle|^2 = \frac{2}{9}.
\]
From this one matrix element follow (Doc.~293, 367, 368):

\begin{enumerate}
\item \textbf{The Koide phase $\theta = 2/9$} \BM: the same number, written
  as an angle in the Koide parametrisation, gives the three charged lepton
  masses without free parameters to $10^{-3}$\,\% (Doc.~292, 293, 367).
\item \textbf{The redistribution weights} $p_0 = 2/9$, $p_1 = (2+3\phi)/9$,
  $p_2 = (5-3\phi)/9$ with $p_0+p_1+p_2 = 1$ \BM.
\item \textbf{The $\phi$ skeleton of the weight ratios} \BM:
  \[
    \frac{p_1}{p_2} = \phi^8, \qquad \frac{p_0}{p_2} = 2\phi^4.
  \]
  Both relations are exact algebraic identities following from
  $p_1=(2+3\phi)/9$ and $p_2=(5-3\phi)/9$ (Doc.~368).
\item \textbf{The link 4/3 = 2Q} (Doc.~371): the Koide quotient $Q = 2/3$ and
  the three-dimensionality factor $4/3$ of the electromagnetic mass
  (Abraham 1902) satisfy $4/3 = 2Q$ exactly.
\end{enumerate}

The mechanism thus carries four independent structural results from a
single matrix-element value $2/9$.

\section{The Spectrum of the Matrix Elements \BM}
\label{sec:spektrum}

Before asking which constant the mechanism might carry further, the reverse
question must be answered: which numbers can it produce at all? This is a
finite calculation. The group $A_5$ has 60 elements; in the
three-dimensional representation, expressed in the $\mathbb Z_3$ mode basis
$v_0,v_1,v_2$, every element has nine matrix elements. The check script
generates all 60 elements from $R_5$ and $C_3$ and collects the squared
moduli. Exactly eight values occur:
\begin{align*}
  &\Bigl\{\,0,\ \tfrac{5-3\phi}{9},\ \tfrac19,\ \tfrac29,\ \tfrac49,\
  \tfrac59,\ \tfrac{2+3\phi}{9},\ 1\,\Bigr\} \\
  &\quad = \{0,\ 0.0162,\ 0.1111,\ 0.2222,\ 0.4444,\ 0.5556,\ 0.7616,\ 1\}.
\end{align*}
All denominators are $9 = |\mathbb Z_3|^2$; the two irrational values are the
weights $p_2$, $p_1$ from Doc.~293/368. The set is not closed under
$x\mapsto 1-x$ ($1-1/9 = 8/9$ is missing), but the subset $\{4/9,5/9\}$ and
the pair $\{p_1, p_0+p_2\}$ are.

\begin{quote}
\textbf{Reach theorem.} Every constant that arises as the squared modulus of
a single $A_5$ matrix element in the $\mathbb Z_3$ basis lies in this
eight-element set. What does not lie in it cannot be delivered by the
mechanism in this form \BM.
\end{quote}

\section{Test Against the Mixing Parameters of the Standard Model}
\label{sec:pruefung}

The eight-element set is held against the dimensionless mixing parameters of
the Standard Model (PDG 2024, NuFIT~6.0). Hits and misses are both recorded.

\subsection{Weinberg angle: on-shell \texorpdfstring{$=2/9$}{= 2/9} \KM}
\label{sec:weinberg}

The Weinberg angle has two common definitions. The $\overline{\rm MS}$ value
$\sin^2\hat\theta_W(M_Z) = 0.2312$ depends on the renormalisation scheme. The
on-shell definition
\[
  \sin^2\theta_W \equiv 1 - \frac{M_W^2}{M_Z^2}
\]
by contrast is fixed by pole masses. Since the FFGFT ground relation
$\tilde T\cdot m = 1$ yields pole masses (Doc.~333), the on-shell scheme is
the natural one for FFGFT.

\begin{center}
\small\begin{tabular}{@{}lcccc@{}}
\toprule
Input & $M_W$ [GeV] & $1-M_W^2/M_Z^2$ & to $2/9$ & Sign. \\
\midrule
World average 2024 (without CDF) & $80.3692$ & $0.22321$ & $0.44\,\%$ & $3.8\sigma$ \\
CDF 2022 & $80.4335$ & $0.22197$ & $0.12\,\%$ & $1.4\sigma$ \\
PDG on-shell direct & --- & $0.22305$ & $0.37\,\%$ & $3.6\sigma$ \\
$\overline{\rm MS}$ (comparison) & --- & $0.23121$ & $3.9\,\%$ & --- \\
\bottomrule
\end{tabular}
\end{center}

The value $2/9$ lies between the world average and the CDF value; the value
of $M_W$ that gives $\sin^2\theta_W = 2/9$ exactly is
\[
  M_W(2/9) = M_Z\sqrt{7/9} = 80.420\ \text{GeV}.
\]
This is a falsifiable prediction. It lies $3.8\sigma$ from the current world
average and $1.4\sigma$ from the CDF value; the $M_W$ measurements are at
present not mutually consistent. The classification is \KM~at $0.4\,\%$,
with the tension stated.

This replaces the previous corpus status: Doc.~336 named as the best Galois
candidate $|\mathrm{GF}(3)^*|/|\mathrm{GF}(9)^*| = 1/4$ with $8\,\%$
deviation. The renormalisation-group running from $\xi$ (Doc.~323, 0.2308)
and the on-shell matrix element $2/9$ are two descriptions: the first in the
$\overline{\rm MS}$ scheme from the scale flow, the second in the on-shell
scheme from the lattice geometry. Whether both describe the same geometric
object remains open \SM; that they serve different schemes is settled.

\subsection{Atmospheric mixing angle: \texorpdfstring{$\sin^2\theta_{23}\in\{4/9,5/9\}$}{sin2 theta23 in \{4/9, 5/9\}} \KM}

The eight-element set contains the pair $4/9$, $5/9$ with $4/9+5/9 = 1$. For
the atmospheric neutrino mixing angle, $\sin^2\theta_{23}$ is experimentally
close to $1/2$, with the octant unresolved:

\begin{center}
\begin{tabular}{@{}lccc@{}}
\toprule
Fit & $\sin^2\theta_{23}$ & nearest candidate & significance \\
\midrule
PDG 2024 (NO) & $0.558\pm0.020$ & $5/9 = 0.5556$ & $0.1\sigma$ \\
NuFIT 6.0 (IO) & $0.550\pm0.015$ & $5/9$ & $0.4\sigma$ \\
NuFIT 6.0 (NO) & $0.470\pm0.015$ & $4/9 = 0.4444$ & $1.7\sigma$ \\
\bottomrule
\end{tabular}
\end{center}

In the matrix-element picture the octant ambiguity of the data is the
ambiguity $4/9\leftrightarrow 5/9$. The mechanism thus makes a prediction
that is independent of the octant:
\[
  \Bigl|\sin^2\theta_{23} - \tfrac12\Bigr| = \tfrac1{18} = 0.0556 .
\]
Maximal mixing ($\sin^2\theta_{23} = 1/2$) is excluded. This can be tested
with DUNE and Hyper-Kamiokande. Classification \KM~for the value, \SM~for
the assignment to the neutrino sector (the neutrino masses themselves come
from $\mathrm{GF}(27)^*$ in Doc.~340, not from $A_5$; that the angle
$\theta_{23}$ is nevertheless an $A_5$ matrix element would be a second,
independent $A_5$ access to the lepton sector).

\subsection{Relation to the Galois reading in Doc.~340}
\label{sec:340}

Doc.~340 reads two mixing angles off the orbit elements of
$\mathrm{GF}(27)^*\cong\mathbb Z_2\times\mathbb Z_{13}$:
$\sin^2\theta_{12}\approx\cos^2(4\pi/13) = 0.323$ and
$\sin^2\theta_{23}\approx\cos^2(10\pi/13) = 0.560$. For $\theta_{23}$ two
values therefore stand side by side, $0.560$ and $5/9 = 0.556$, both of which
hit the PDG value to within $0.4\,\%$. This needs to be clarified.

\paragraph{The two values are not identical \BM.}
$\cos(2\pi/13)$ has a minimal polynomial of degree $\varphi(13)/2 = 6$ over
$\mathbb Q$; consequently $\cos^2(10\pi/13)$ is irrational, whereas $5/9$ is
rational. The closeness of $0.85\,\%$ is numerical, not algebraic: $5/9$ is
merely the best rational approximation to $\cos^2(10\pi/13)$ with denominator
$\le 13$.

\paragraph{The Galois reading is an order-of-magnitude statement.}
The set $\{\cos^2(2\pi k/13),\ 1-\cos^2(2\pi k/13)\mid k=1..6\}$ contains
twelve values between $0.01$ and $0.99$. Held against seven dimensionless
parameters of the Standard Model, for six of them an element is found within
$6.6\,\%$ --- the set is too dense to discriminate at this tolerance. The
eight-element set of §\ref{sec:spektrum}, by contrast, discriminates: three
hits below $1.3\,\%$, the rest above $5\,\%$. The $\cos^2$ readings in
Doc.~340 are therefore to be understood as orders of magnitude, not as
derivations; Doc.~340 itself labels them as ``read off''.

\paragraph{Division of labour in the neutrino sector.}
This determines which structure carries which parameter:
\begin{center}
\small\begin{tabular}{@{}llll@{}}
\toprule
Quantity & Structure & Source & Status \\
\midrule
$\Delta m^2_{\rm atm},\ \Delta m^2_{\rm sol}$ & $\mathrm{GF}(27)^*$ orbits & Doc.~340 & \KM~(1\,\%, 0.5\,\%) \\
$\nu_3$ Dirac, $\nu_{1,2}$ Majorana & quadratic character mod 13 & Doc.~343 & \BM \\
$\sin^2\theta_{13} = (25\xi)^{2/3}$ & $\xi\cdot|A_5|$ ($25\xi = 1/300$) & Doc.~340 & \KM~(1.4\,\%) \\
$\sin^2\theta_{23}\in\{4/9,5/9\}$ & $A_5$ matrix element & Doc.~372 & \KM~(0.4\,\%) \\
$\sin^2\theta_{12}$ & only $\cos^2(4\pi/13)$, order of magnitude & Doc.~340 & \SM~(5\,\%) \\
$\delta_{CP}$ & orbit-3 element $k=10$ & Doc.~340 & \SM \\
\bottomrule
\end{tabular}
\end{center}
The neutrino sector is carried by two structures: $\mathrm{GF}(27)^*$ for the
masses and the Dirac/Majorana character, $A_5$ (alone or with $\xi$) for
$\theta_{23}$ and $\theta_{13}$. Both sit on the same lattice (Doc.~364:
$A_5\subset\mathrm{Aut}(D_4)$ only via the phase channel $\mathrm{GF}(81)$;
$\mathrm{GF}(27)$ as the ordering structure, Doc.~341). Whether $\theta_{12}$
is likewise an $A_5$ object that only becomes visible through a product or
another basis, or whether it remains with $\mathbb Z_{13}$ alone, is the
remaining open question of the sector \SM.

\subsection{Cabibbo angle \SM}

$\lambda = \sin\theta_C = 0.2250\pm0.0007$ lies $1.2\,\%$ ($4.1\sigma$) from
$2/9$. The long-known numerical closeness
$\sin\theta_C\approx\sin^2\theta_W^{\rm on\text{-}shell}$ is read here as
the closeness of both quantities to $2/9$. This is not an identification:
Doc.~041 derives $|V_{us}| = \sqrt{m_d/m_s}\cdot f_{\rm Cab}$ from quark
masses, and the quark coefficients are QCD-calibrated (Doc.~289). Whether
$\lambda$ is a matrix element of another subgroup of $\mathrm{Aut}(D_4)$ is
open \SM.

\subsection{Negative findings \XM}

Three parameters do not lie in the eight-element set and cannot be reached
by a single $A_5$ matrix element:
\begin{center}
\begin{tabular}{@{}lccc@{}}
\toprule
Parameter & Value & nearest element & distance \\
\midrule
$\sin^2\theta_{12}$ (solar) & $0.307$ & $2/9$ & $28\,\%$ \\
$\sin^2\theta_{13}$ (reactor) & $0.0222$ & $(5-3\phi)/9$ & $27\,\%$ \\
$\alpha_s(M_Z)$ & $0.118$ & $1/9$ & $5.8\,\%$ \\
\bottomrule
\end{tabular}
\end{center}
The reach theorem of §\ref{sec:spektrum} makes this a sharp statement: these
three quantities require either another group, a product or difference of
matrix elements, or a different mechanism. For $\alpha_s$ the corpus has the
Galois running with $\xi$; for $\theta_{12}$ and $\theta_{13}$ nothing is
available in the matrix-element picture.

\subsection{Quarks: compact description present, unrolling incompatible}
\label{sec:quarks}

In the corpus the charged leptons have two descriptions that meet: compact,
the Yukawa ladder $m = r\,\xi^p v$ with $r_e = 4/3$ (Doc.~006), and unrolled
($T^4\to\mathbb R^3\times S^1$, Doc.~370), the $A_5$ mechanism of this
document. Both arrive at $Q = 2/3$, $r = \sqrt2$ --- the ``special point''
of Doc.~289.

The quarks have only the compact description, and it is more developed than
the comparison with unrolling suggests. The exponents $p_i$ are derived from
the torus topology (Doc.~189: $p = d_{\rm akt}/3$, one active dimension fewer
per step) --- up/down $3/2$, strange $1$, charm $2/3$, bottom $1/2$ as
transition exponent, top $-1/3$ as inverse winding \KM. The $\xi$ scale
correction and the cumulative factor $K_{\rm frak} = 1-100\xi$ from the scale
flow Planck$\to$electroweak (Doc.~133) weight the ladder; the extended
fractal method thereby reaches $\sim1.2\,\%$ without free parameters
(Doc.~005). What is missing are the prefactors $r_i$: Doc.~189 sets
$r_i = n_\phi^2/(n_\theta\,n_c)$ with $n_c = 3$ for quarks and explicitly
names the complete derivation as outstanding; the $r_i$ tabulated in
Doc.~006 are back-calculated from the measured masses, but carry exclusively
prime factors from the Galois grid $\{2,3,5,7,11,13\}$ \BM~(Doc.~358,
Theorem~B). The gap is therefore precise: $p_i$ derived, $r_i$ postulated
\SM.

The unrolling cannot be transferred, and Doc.~289 shows why: the quark
triplets do not sit at the special point. With PDG-like values,
$Q_{\rm up} = 0.849$ ($r = 1.76$) and $Q_{\rm down} = 0.730$ ($r = 1.54$).
Scale dependence does not change this: running all six masses to a common
scale ($\overline{\rm MS}$ at $M_Z$, $m_t$ or roughly $M_{\rm GUT}$),
$Q_{\rm up}$ stays between $0.85$ and $0.89$ and $Q_{\rm down}$ between $0.72$
and $0.75$ --- never close to $2/3$ \KM. The $A_5$ mechanism, however,
produces \emph{exactly} $Q = 2/3$ --- that is its signature
($\theta = |A|^2 = 2/9 \Rightarrow Q = 2/3$, Doc.~367), and the eight-element
set contains no value that yields $Q = 0.85$ or $0.73$. The quark masses are
incompatible with the lepton unrolling mechanism \XM, not merely unworked.

Three structural reasons, each documented in the corpus: quarks carry
colour, and the $\mathbb Z_3$ that counts the generations for leptons is at
the same time the colour $\mathbb Z_3$ for quarks (Doc.~321) --- the mode
space would be $3\times3$, not~3. Quarks come as isospin doublets, six
instead of three modes per sector. And $\tilde T\cdot m = 1$ gives pole
masses, which are not defined for $u,d,s$ (Doc.~333). The open question is
therefore not ``how does one unroll quarks like leptons'' but: which other
subgroup of $\mathrm{Aut}(D_4)$ or which tensor product
$A_5\otimes\mathbb Z_3^{\rm colour}$ has a matrix-element spectrum that
contains $Q\approx0.85$ and $0.73$. Doc.~358 names the field
$\mathrm{GF}(729)$ as the entry point for the top coefficient (prime 7,
$k=6$) \SM. That is a different programme, not an appendix to lepton
unrolling.

\subsection{Gauge bosons and Higgs}
\label{sec:bosonen}

\paragraph{What is derived.}
The electroweak scale $v$ follows from the lattice structure to $0.3\,\%$
(Doc.~006, R104: $248.3$\,GeV without, $245.6$\,GeV with $K_{\rm frak}$)
\KM. The ratio $M_W/M_Z$ follows from §\ref{sec:weinberg}:
$M_W^2/M_Z^2 = 1 - 2/9 = 7/9$ \KM. The fine-structure constant
$1/\alpha = 3700/27$ follows from Galois orders (Doc.~338) \KM. Thus $v$,
$\alpha$ and $\sin^2\theta_W$ are available; the absolute masses
$M_W = gv/2$, $M_Z = M_W/\cos\theta_W$ (Doc.~320) follow from them via the
standard electroweak relations up to radiative corrections ($\Delta r$, of
the order of a per cent), which the corpus does not compute. Photon and the
eight gluons are massless: through the unbroken $U(1)_{\rm EM}$ and
$SU(3)_c$ (Doc.~320) \BM, where $SU(3)_c$ follows from $\mathbb Z_3$
triality (Doc.~321) and the eight gluons from the eight three-element orbits
of $\mathrm{GF}(27)^*$ (Doc.~339) \BM.

\paragraph{Where the Higgs stands on the ladder.}
The $p$ ladder of Doc.~189 places the Higgs scale at $p = 0$
($d_{\rm akt} = 0$): $m_h = r_h\,v$ without $\xi$ suppression, with
$r_h = m_h/v = 0.508$. This correctly places the Higgs as the only particle
on the $v$ scale but provides no prefactor --- the same gap as for the quark
$r_i$.

\paragraph{What is not derived: the Higgs mass \XM.}
The corpus contains two formulas for $m_h$, and neither survives
recalculation. Doc.~041 gives $m_h = v\,\xi^{1/4}$ with the numerical value
$125$\,GeV; in fact $v\,\xi^{1/4} = 26.5$\,GeV. The same table gives
$\Lambda_{\rm QCD} = v\,\xi^{1/3} = 200$\,MeV; in fact $12.6$\,GeV. Doc.~005
gives $m_H = m_t\,\phi\,(1+\xi D_f) = 124.95$\,GeV; since
$m_t\phi = 279.4$\,GeV, this requires $\xi D_f = -0.553$, i.e.\
$D_f\approx-4100$ --- a fit parameter, not a derivation. These two formulas
are recorded as R122 in Doc.~190. Viable candidates are only provided by
Doc.~373 (Galois chain and trace rule). The eight-element set yields nothing
either: $m_h/v = 0.508$ and $m_h^2/v^2 = 0.258$ do not lie in it \XM.

\paragraph{Systematic search.}
Instead of guessing individual formulas, 17 dimensionless ratios of $m_h$,
$M_W$, $M_Z$, $m_t$, $v$ were held against a library of 276 FFGFT-related
numbers (fractions with denominator $\le 12$, powers of $\phi$, roots, Galois
orders, $K_{\rm frak}$, the eight-element set). At $0.5\,\%$ tolerance twelve
hits are found, at $0.2\,\%$ three; the random expectation is about fourteen
and six respectively. There is therefore \emph{no} signal above the
look-elsewhere background \XM. The obvious reading $m_h = v/2$ lies at
$1.7\,\%$ and is not distinguished.

A single candidate has internal structure and is therefore recorded,
explicitly as selected after the fact: $m_h/M_Z = 1.3730 \pm 0.0012$ and
$m_t/m_h = 1.3784 \pm 0.0026$ both lie at $11/8 = 1.375$ (at $1.7\sigma$ and
$1.3\sigma$ respectively), hence
\[
  M_Z : m_h : m_t \;\approx\; 1 : \tfrac{11}{8} : \bigl(\tfrac{11}{8}\bigr)^2,
  \qquad m_h \approx \sqrt{m_t\,M_Z} \quad (0.20\,\%,\ 1.6\sigma).
\]
The numbers $11 = \max(\text{orbit}_4)$ and $8 = |\mathrm{GF}(9)^*|$ are both
Galois quantities of the corpus (Doc.~340, 336). The Higgs would then sit as
the geometric mean between $Z$ and top --- in the language of
$\tilde T\cdot m = 1$: its winding number is the mean of the two. Status
\SM, with the proviso that one candidate out of about 5000 tests is not
evidence but a hypothesis that would need an independent derivation.


\section{Hadrons: Bound States on the Lattice}
\label{sec:hadronen}

The leptons are unrolled primary modes of the $T^4/D_4$ lattice --- free,
asymptotic particles whose masses follow directly from
$\tilde T\cdot m = 1$. Hadrons are not. There is no free quark; what is
measured is the system energy of the bound state. Three consequences can be
tested numerically.

\paragraph{Koide quotient of the hadron multiplets \BM.}
For an isospin multiplet of dimension $n$ with nearly degenerate masses
$m_i = m(1+\delta_i)$, $|\delta_i|\ll1$, it holds algebraically exactly that
\[
  Q \cdot n = 1, \qquad Q = \frac{\sum m_i}{\bigl(\sum\sqrt{m_i}\bigr)^2}.
\]
Numerically: $Q(p,n) = 1/2$, $Q(\Sigma^+,\Sigma^0,\Sigma^-) = 1/3$,
$Q(\Delta^{++},\Delta^+,\Delta^0,\Delta^-) = 1/4$, $Q(\Omega^-) = 1$, all to
$< 0.01\,\%$ \BM. This is an algebraic identity for degenerate masses, not a
new result --- but it sharpens the contrast: the leptons break exactly this
relation. For $n = 3$ generations the degeneracy rule would require
$Q_{\rm Lep} = 1/3$; in fact $Q_{\rm Lep} = 2/3$. The factor~2 is the
signature of the $A_5$ mechanism: the icosahedron lifts the degeneracy
specifically such that $Q = 2\cdot(1/3) = 2/3$. Hadrons remain at the
degenerate relation because their $SU(2)$ isospin symmetry is almost exactly
conserved.

\paragraph{Gell-Mann--Okubo splitting consistent with $m_s$ \KM.}
The baryon decuplet masses show a uniform spacing:
$\Delta m(\Sigma^*{-}\Delta) = 151.7\,\text{MeV}$,
$\Delta m(\Xi^*{-}\Sigma^*) = 148.1\,\text{MeV}$,
$\Delta m(\Omega{-}\Xi^*) = 140.4\,\text{MeV}$, mean $146.7\,\text{MeV}$.
This is the Gell-Mann--Okubo rule; the splitting corresponds to the effective
strange quark mass. From the Galois coefficient $r_s$ (Doc.~006) follows
$m_s^{(\rm FFGFT)} = r_s\,\xi\,v \approx 93\,\text{MeV}$. The agreement is
consistent to a factor of~1.6 \KM, but not a precision statement --- the
difference lies in the QCD dynamics, not in the lattice geometry.

\paragraph{$\Lambda_{\rm QCD}$ not derivable from the $\xi$ ladder \XM.}
$\Lambda_{\rm QCD} \approx 210\,\text{MeV}$ (PDG, $n_f = 5$) could be written
as $v\,\xi^p$ for $p \approx 0.79$ --- not a step of the ladder $p = n/3$ or
$n/2$. The route via the hadron radius ($\omega_{\min} = 1/R_{\rm had}$,
Doc.~319) is also circular as long as $R$ does not follow from $\xi$.
$\Lambda_{\rm QCD}$ is thus the sharpest open question of the hadron sector:
it lies neither in the eight-element set (§\ref{sec:spektrum}), nor on the
$\xi$ ladder, nor in the Galois grid of Doc.~358.

\paragraph{Gluons as standing waves.}
Doc.~319 describes gluons as standing waves in the closed proton torus:
massless, pure field energy ($E = pc$, $d\tau = 0$), topologically confined
by the $\mathbb Z_3/SU(3)$ boundary condition (eight linking states of the
three flux strands = eight generators of $SU(3)_c$, Doc.~321 \BM). In this
picture confinement is not a dynamical mechanism but a topological
impossibility of escape: the closed torus admits only integer wave numbers;
there is no free propagation to the outside. In this language
$\Lambda_{\rm QCD}$ is the minimum frequency on the torus,
$\omega_{\min} = 1/R_{\rm Proton}$ (Doc.~319 \SM).

\paragraph{$R_{\rm Proton}$ structurally classified \SM.}
Doc.~319 gives $\Lambda_{\rm QCD} \approx \omega_{\min} = \hbar c/R_{\rm Proton}$
\SM. The FFGFT quantities allow a structural classification: the geometric
mean of the classical electron radius $r_e = \alpha\,\hbar c/m_e$ and the
proton Compton wavelength $\lambda_p = \hbar c/m_p$ is
\[
  R_{\rm geom} = \sqrt{r_e\,\lambda_p}
  = \frac{\hbar c\,\sqrt{\alpha}}{\sqrt{m_e\,m_p}}
  = 0.770\,\text{fm},
\]
$8.5\,\%$ below the measured charge radius $0.841\,\text{fm}$. All three
factors --- $\alpha$ from Galois structure (Doc.~338), $m_e$ from the $A_5$
mechanism, $m_p$ from the $\xi$ ladder (Doc.~189) --- are available in FFGFT.
The missing factor $R_{\rm PDG}/R_{\rm geom} \approx 1.09$ lies in the
quark--gluon binding dynamics, which the programme does not yet capture
formally.

\paragraph{Conclusion.}
Hadrons are bound states on the same lattice as leptons --- the same $\xi$,
the same ladder for $p_i$, the same Galois grid for $r_i$. Confinement is a
topological boundary condition, not a separate mechanism. $R_{\rm Proton}$
is structurally classified to $8.5\,\%$; the residual lies in the binding
dynamics. This is the honest limit \SM.

\section{The More General Principle}
\label{sec:prinzip}

The four results of §\ref{sec:ertrag} suggest a more general thesis:

\begin{quote}
\textbf{The matrix-element principle}: every dimensionless physical constant
that arises from the overlap of two lattice modes under a finite group
rotation can be read as the squared modulus of a matrix element
$|\langle v_i|g|v_j\rangle|^2$ --- where $g$ is a group element, $v_i, v_j$
are orthonormal modes of the $T^4/D_4$ lattice, and the result has no free
parameters.
\end{quote}

For the lepton masses this principle is fully carried out with
$g = R_5 \in A_5$, $i=0$, $j=1$ \BM. For other sectors it is an open
hypothesis \SM. The sharpest test question is: is there a finite subgroup of
$\mathrm{Aut}(D_4)$ and two lattice modes whose squared matrix-element
modulus reproduces the Weinberg angle or the Cabibbo angle without free
parameters?

That $|A_5|= 60 = 4 \cdot 3 \cdot 5$ unites the three defining symmetry
numbers of the lepton structure, and that $A_5$ is the only simple finite
subgroup of $SO(3)$ with five-fold elements, make this hypothesis
structurally attractive.

\section{Result}

\begin{enumerate}
\item \textbf{Reach theorem} \BM: in the $\mathbb Z_3$ basis $A_5$ generates
  exactly eight squared matrix-element moduli
  $\{0, (5-3\phi)/9, 1/9, 2/9, 4/9, 5/9, (2+3\phi)/9, 1\}$.
\item \textbf{Lepton masses}: from $2/9$ follow the Koide phase, the weights,
  the $\phi$ skeleton ($p_1/p_2 = \phi^8$, $p_0/p_2 = 2\phi^4$) and
  $4/3 = 2Q$ \BM.
\item \textbf{Weinberg angle} \KM: on-shell $1-M_W^2/M_Z^2 = 2/9$ to
  $0.4\,\%$; prediction $M_W = M_Z\sqrt{7/9} = 80.420$\,GeV, between the
  world average and CDF. Replaces the candidate $1/4$ (8\,\%) of Doc.~336.
  The $\overline{\rm MS}$ result 0.2308 of Doc.~323 remains as the scale-flow
  description.
\item \textbf{Atmospheric mixing} \KM: $\sin^2\theta_{23}\in\{4/9,5/9\}$,
  prediction $|\sin^2\theta_{23}-1/2| = 1/18$; maximal mixing excluded. The
  PDG value $0.558$ hits $5/9$ at $0.1\sigma$.
\item \textbf{Comparison with Doc.~340} \BM: $\cos^2(10\pi/13)\ne 5/9$
  (degree 6 versus rational); the $\cos^2(2\pi k/13)$ readings are orders of
  magnitude. Division of labour: $\mathrm{GF}(27)^*$ carries masses and
  Dirac/Majorana, $A_5$ carries $\theta_{23}$ and (with $\xi$) $\theta_{13}$;
  $\theta_{12}$ remains open \SM.
\item \textbf{Quarks} \XM: compact ladder with derived $p_i$ (Doc.~189) and
  $\xi$/$K_{\rm frak}$ weighting (Doc.~133, 005), $\sim1.2\,\%$ without free
  parameters; $r_i$ postulated, in the Galois grid \BM.
  $Q_{\rm up}\approx0.85$--$0.89$, $Q_{\rm down}\approx0.72$--$0.75$ at all
  scales: incompatible with $A_5$ unrolling ($Q = 2/3$ exactly).
\item \textbf{Bosons}: $v$, $\alpha$, $\sin^2\theta_W$ available \KM;
  photon/gluons massless \BM; the Higgs formulas in Doc.~041 and 005 do not
  survive recalculation \XM~(R122); candidates in Doc.~373.
\item \textbf{Cabibbo} \SM: $\lambda$ within $1.2\,\%$ of $2/9$, not
  identified; Doc.~041 has a different derivation.
\item \textbf{Negative} \XM: $\sin^2\theta_{12}$, $\sin^2\theta_{13}$,
  $\alpha_s$ do not lie in the eight-element set.
\item The general matrix-element principle (§\ref{sec:prinzip}) thus has
  three confirmed instances ($2/9$ lepton, $2/9$ Weinberg, $\{4/9,5/9\}$
  atmospheric) and three excluded ones; the next question is whether another
  subgroup of $\mathrm{Aut}(D_4)$ carries the excluded values.
\end{enumerate}

\medskip
\noindent\textbf{Check script:}\newline\texttt{2/python/Dok372\_Skripte/pruef\_372\_matrixelement.py}

\medskip
\noindent\textbf{Register (Doc.~190):} R122 (Higgs mass in Doc.~041 and 005).

\begin{thebibliography}{9}
\bibitem{dok292} J.~Pascher, \emph{Doc.~292 --- Koide formula and FFGFT} (2025).
\bibitem{dok293} J.~Pascher, \emph{Doc.~293 --- Icosahedral matrix element} (2025).
\bibitem{dok323} J.~Pascher, \emph{Doc.~323 --- Renormalisation group and Weinberg angle} (2025).
\bibitem{dok340} J.~Pascher, \emph{Doc.~340 --- Neutrino mass hierarchy from GF(27)*} (2026).
\bibitem{dok367} J.~Pascher, \emph{Doc.~367 --- Why $\theta = 2/9$} (2026).
\bibitem{dok368} J.~Pascher, \emph{Doc.~368 --- From $p_0,p_1,p_2$ to lepton masses} (2026).
\bibitem{dok371} J.~Pascher, \emph{Doc.~371 --- The factor 4/3} (2026).
\end{thebibliography}
